% concatenated from ARX_FUFCL_140624.tex
\documentclass[12pt]{article}
\usepackage{amsmath, amsthm}\usepackage{enumerate}
\usepackage{amssymb}\usepackage{bbold}
\usepackage{color}
\usepackage[explicit]{titlesec}
\newtheorem{thm}{Theorem}
\newtheorem{prop}{Proposition}
\newtheorem{lem}{Lemma}
\newtheorem{exam}{Example}
\newtheorem{cor}{Corollary}
\newtheorem{rem}{Remark}
\newtheorem{defi}{Definition}
\newtheorem{ques}{Question}
\DeclareMathOperator{\ruc}{\xrightarrow[]{r}}
\DeclareMathOperator{\oc}{\xrightarrow[]{o}}
\DeclareMathOperator{\ic}{\xrightarrow[]{i}}
\DeclareMathOperator{\mc}{\xrightarrow[]{\mathcal{M}}}
\setlength{\parindent}{0in}
\renewcommand{\le}{\leqslant}
\renewcommand{\ge}{\geqslant}
\renewcommand{\mid}{\::\:}
\newcommand{\R}{\mathbb{R}}
\newcommand{\N}{\mathbb{N}}
\newcommand{\veps}{\varepsilon}
\newcommand{\blankbox}[2]{
\parbox{\columnwidth}{\centering
\fbox{\raisebox{0pt}[#2]{\hspace{#1}}}}}
\renewcommand{\baselinestretch}{0.98}
\makeatletter
\renewcommand{\subsection}{\@startsection{subsection}{1}
{0pt}{3.25ex plus 1ex minus.2ex}{-1em}{\normalfont\normalsize\bf}}
\begin{document}
\title{Free uniformly complete vector lattices}
\date{\empty}
\maketitle
\author{\centering{Eduard Emelyanov and Svetlana Gorokhova\\ 
}}
\maketitle

\medskip
\medskip
\medskip

{\bf{Keywords:} \rm relatively uniform convergence, free vector lattice, 
free uniformly complete vector lattice, free Banach lattice}

\medskip

{\bf  Mathematics Subject Classification:} {\normalsize  46A40, 46B42}

%\medskip
%
%{\bf  Abstract:} {\normalsize We define a free uniformly complete 
%vector lattice $\text{\rm FUCVL}(A)$ over a non-empty set $A$ and give its representation 
%as a sublattice of the space $H(\mathbb{R}^A)$ of continuous in the product topology positively 
%homogeneous functions on $\R^A$.}


\section{Introduction and preliminaries}

\subsection{Free vector lattice over a set.}

We shall abbreviate a vector lattice (sublattice) by VL (resp., VSL).
Following the approach of de Pagter and Wickstead (cf. \cite[Definition 3.1]{PW}),
a {\em free vector lattice over} a non-empty set $A$ (briefly, $\text{\rm FVL}(A)$) 
is a pair $(F,i)$, where $F$ is a VL and $i:A\to F$ is a map such that, 
for any VL $E$ and for any map $q :A\to E$, there exists 
a unique Riesz homomorphism $T: F \to E$ satisfying 
$q=T \circ i$. If $(F,i)$ is a $\text{\rm FVL}(A)$ then $F$ is generated by $i(A)$ as a VL. 
The existence of a $\text{\rm FVL}(A)$ over $A$ is the long 
established fact going back to Birkhoff \cite{Birk}. A concrete representation 
of $\text{\rm FVL}(A)$ as a VL of real-valued functions was 
given by Weinberg \cite{Wein} and Baker \cite{Baker} 
(cf. also \cite{Blei}). Notice that $\text{\rm FVL}(A)$ is always Archimedean.

\smallskip
A {\em free Banach lattice} (briefly, $\text{\rm FBL}(A)$) {\em over 
a non-empty set} $A$ was introduced and investigated in \cite{PW} 
(see also \cite{Tr} for a simple alternative approach to a free 
lattice norm on $\text{\rm FVL}(A)$). In the present note, we define 
a {\em free uniformly complete vector lattice over} 
$A$ and give some of its representations. For further unexplained 
notation and terminology, we refer to \cite{Ku,PW}.

\subsection{Uniformly complete vector lattices.}
A net $x_\alpha$ in a VL $U$ {\em relatively uniformly converges} 
(briefly, {\em $\text{\rm r}$-converges}) 
to $x\in U$ if there exists $u \in U$ ({\em regulator of convergence}) 
such that, for each $k\in\N$, there exists $\alpha_k$ 
with $|x_\alpha-x|\le\frac{1}{k}|u|$ 
for all $\alpha\ge\alpha_k$ 
(in this case, we write $x_\alpha\ruc x(u)$). 
A net $x_\alpha$ in a VL $U$ is {\em $\text{\rm r}$-Cauchy} with a regulator $u\in U$
(briefly, $\text{\rm r}(u)$-{\em Cauchy}) if, for each $k\in\N$, there exists $\alpha_k$
with $|x_{\alpha'}-x_{\alpha''}|\le\frac{1}{k}|u|$ for all $\alpha',\alpha''\ge\alpha_k$.

\smallskip
A VSL $V$ of a VL $U$ is called \text{\rm r}-{\em complete in} $U$ if
$U$ is Archimedean and, for every $u\in U$ and every $\text{\rm r}(u)$-Cauchy
net $x_\alpha$ in $V$, there exist $x\in V$ and $w\in U$ 
with $x_\alpha\ruc x(w)$ (such an $x$ is unique since $V$ is Archimedean). 
For example, VSL $c_{00}$ of $c_0$ is not \text{\rm r}-complete in $c_0$ as 
$c_{00}\ni\sum_{i=1}^{n}\frac{1}{k^2}e_i\ruc
\sum_{i=1}^\infty\frac{1}{k^2}e_i\notin c_{00}$
with regulator $\sum_{i=1}^\infty\frac{1}{k}e_i\in c_0$, 
yet $c_{00}$ is \text{\rm r}-complete in itself.

\smallskip
We include the following, certainly well known fact, for which we did not 
found an appropriate reference.

\begin{lem}[]\label{L1-0}
Let $x_\alpha$ be an $\text{\rm r}(u)$-Cauchy
net in an Archi\-me\-dean \text{\rm VL} $U$, where $u\in U$.
If $x_\alpha\ruc x(w)$ with $x,w\in U$, then $x_\alpha\ruc x(u)$.
\end{lem}

\begin{proof}
For each $l\in\N$ take an $\alpha(l)$ such that $|x_\alpha-x|\le\frac{1}{l}|w|$ 
when $\alpha\ge\alpha(l)$. 
Since $x_\alpha$ is $\text{\rm r}(u)$-Cauchy,
for each $k\in\N$, there exists $\alpha_k$ with 
$|x_{\alpha'}-x_{\alpha''}|\le\frac{1}{k}|u|$
for $\alpha',\alpha''\ge\alpha_k$. Now, for every $l,k\in\N$, 
we pick $\alpha(k,l)\ge\alpha_k,\alpha(l)$. 
Then, for each $\alpha\ge\alpha_k$, 
$$
   |x_\alpha-x|\le|x_\alpha-x_{\alpha(k,l)}|+
   |x_{\alpha(k,l)}-x|\le\frac{1}{k}|u|+\frac{1}{l}|w|.
   \eqno(1)
$$
Since $l$ is arbitrary in (1) and $U$ is Archimedean then
$|x_\alpha-x|\le\frac{1}{k}|u|$ for $\alpha\ge\alpha_k$,
and hence $x_\alpha\ruc x(u)$.
\end{proof}

\begin{cor}[]\label{same reg}
Let $V$ be a \text{\rm VSL} of an Archimedean \text{\rm VL} $U$.
The following conditions are equivalent.
\begin{enumerate}[$i)$]
\item
$V$ is $\text{\rm r}$-complete in $U$.
\item
For every $u\in U$ and every $\text{\rm r}(u)$-Cauchy net $x_\alpha$ in $V$,
there exists $x\in V$ with $x_\alpha\ruc x(u)$.
\end{enumerate}
\end{cor}


\smallskip
As the \text{\rm r}-convergence is sequential
(cf., e.g. \cite[Proposition 2.4]{AEG}), 
we can restrict ourselves to \text{\rm r}-convergent/\text{\rm r}-Cauchy sequences. 
In particular, a VSL $V$ of an Archimedean VL $U$ is \text{\rm r}-complete in $U$
iff, for every $u\in U$ and every $\text{\rm r}(u)$-Cauchy sequence $x_n$ in $V$,
there exists $x\in V$ such that $x_n\ruc x(u)$. 

\smallskip
An \text{\rm r}-complete in itself VL will be abbreviated by \text{\rm UCVL}.
For example, Dedekind complete VLs and Banach lattices are UCVLs.
Note that each UCVL is Archimedean by definition.

\begin{lem}[]\label{princ-ideal}
Let $U$ be a \text{\rm VL}. The following conditions are equivalent.
\begin{enumerate}[$i)$]
\item
$U$ is a \text{\rm UCVL}.
\item
For every $y\in U$, the principal order ideal $U_y=\bigcup_{n\in\mathbb{N}}[-n|y|,n|y|]$ 
is complete under the norm $\|x\|_y:=\inf\{\lambda\in\mathbb{R}:|x|\le\lambda|y|\}$.
\item
$U_y$ is \text{\rm UCVL} for all $y\in U$.
\end{enumerate}
\end{lem}

\begin{proof}
$i)\Longrightarrow ii)$\
Let $y\in U$. As $U$ is Archimedean, the map $x\to\|x\|_y$ is a norm on $U_y$.
Since $\|x_n-x\|_y\to 0\Longleftrightarrow x_n\xrightarrow[]{r}x(|y|)$ 
for each sequence $x_n$ in $U$ and since $U$ is \text{\rm r}-complete, then $U_y$ 
is $\|\cdot\|_y$-complete.

\smallskip
$ii)\Longrightarrow i)$\
Since the map $x\to\|x\|_y$ is a norm on $U_y$ for every $y\in U$, then $U$ is Archimedean.
Let $u\in U$ and let $x_n$ be an $\text{\rm r}(u)$-Cauchy sequence in $U$.
Then, there exists $n_1$ satisfying $|x_k-x_m|\le|u|$ for all $k,m\ge n_1$.
Therefore, $|x_k-x_m|\le w:=|u|+\sum_{i=1}^{n_1}|x_i|\in U$ for all $k,m\in\mathbb{N}$,
and hence the sequence $x_n$ lies in $U_w$. Since $x_n$ is  $\text{\rm r}(u)$-Cauchy 
then it is  $\text{\rm r}(w)$-Cauchy,
and hence it is $\|\cdot\|_w$-Cauchy in $U_w$. Since $U_w$ is $\|\cdot\|_w$-complete,
there exists $x\in U_w$ with $\|x_n-x\|_w\to 0$, and hence $x_n\ruc x(w)$.
Then $x_n\ruc x(u)$ by Lemma \ref{L1-0}. 

\smallskip
$ii)\Longleftrightarrow iii)$\
It is trivial.
\end{proof}

\begin{defi}[]\label{d00}
A \text{\rm UCVL} $F$ is an \text{\rm r}-{\em completion} 
of a \text{\rm VL} $E$ if there exists a Riesz embedding $i: E\to F$ such that, 
for each \text{\rm UCVL} $G$ and each Riesz homomorphism $T: E\to G$ 
there is unique Riesz homomorphism $S: F\to G$ satisfying $T=S\circ i$.
\end{defi}

Since every UCVL is Archimedean, each VL possessesing an \text{\rm r}-completion
is Archimedean. 

\smallskip
If an $\text{\rm r}$-completion $F$ of a VL $E$
exists, it must be unique up to a Riesz isomorphism. Indeed,
let a UCVL $F_1$ be another $\text{\rm r}$-completion of $E$. 
Then, there are two 
Riesz embeddings $i: E\to F$ and $i_1: E\to F_1$ with two 
Riesz homomorphisms $S: F\to F_1$ and $S_1: F_1\to F$ 
satisfying $i_1=S\circ i$ and $i=S_1\circ i_1$
and hence $i=S_1\circ S\circ i$. 
Since a Riesz homomorphism $H: F\to F$ satisfying $i=H\circ i$ is unique then $S_1\circ S=I_F$,
where $I_F$ is the identity map on $F$.
Similarly, $S\circ S_1=I_{F_1}$. Thus, VLs $F$ and $F_1$ are Riesz isomorphic.

\smallskip
It is long known (e.g., see Veksler's works \cite{Veks1,Veks2}) that the intersection 
of all \text{\rm r}-complete VSLs containing $E$ in the Dedekind completion 
$E^\delta$ of an Archimede\-an VL $E$ is an $\text{\rm r}$-completion of $E$. 
For convenience of the reader, we include details of the construction of $\text{\rm r}$-completion 
in slightly more general setting. 

\subsection{$\text{\rm r}$-Completion of an Archimedean vector lattice.}
Throughout this subsection, $E$ is a \text{\rm VSL} of a \text{\rm UCVL} $U$.

\begin{lem}\label{Rem1}
The intersection $F$ of all \text{\rm UCVL}s in $U$ containing $E$ is a \text{\rm UCVL}.
\end{lem}

\begin{proof}
Suppose $x_n$ is an $\text{\rm r}(x_1)$-Cauchy sequence in $F$. 
Take a UCVL $V$ in $U$ containing $E$. Then $x_n\ruc x(x_1)$ for some $x\in V$ 
by Corollary \ref{same reg}. Since $V$ is arbitrary UCVL in $U$ containing $E$, 
and since $\text{\rm r}$-limit $x$ of the sequence $x_n$ is unique, $x$ belongs 
to all UCVLs in $U$ containing $E$, and hence $x\in F$.
\end{proof}

Define a transfinite sequence $(F_\beta)_{\beta\in{\text{\rm Ord}}}$ of VSLs of $U$ by:
\begin{enumerate}[-]
\item 
$F_1:=E$;
\item
$F_{\beta+1}:=\{x\in U: x_n \ruc x(x_1),\ 
\text{\rm for a sequence}\ x_n\ \text{\rm in}\ F_{\beta}\}$;
\item
$F_\beta:=\bigcup_{\gamma\in{\text{\rm Ord}};\ \gamma<\beta}F_{\gamma}$
for a limit ordinal $\beta$.
\end{enumerate}

Note that $F_{\beta_1}$ is a VSL of $F_{\beta_2}$ if $\beta_1\le\beta_2$.

\begin{lem}[]\label{L1}
The intersection $F$ of all \text{\rm UCVL}s of $U$ containing $E$ satisfies
$F=\bigcup_{\gamma\in{\text{\rm Ord;$\gamma<\omega_1$}}}F_{\gamma}$.
\end{lem}

\begin{proof}
Take an $\text{\rm r}(x_1)$-Cauchy sequence $x_n$ in 
$G:=\bigcup_{\gamma\in{\text{\rm Ord;$\gamma<\omega_1$}}}F_{\gamma}$. 
Since $U$ is a UCVL, $x_n\ruc x(x_1)$ for some $x\in U$ by Lemma \ref{L1-0}. 
Pick an ordinal $\beta<\omega_1$ with $x_n\in F_\beta$ for all $n$.
Then $x\in F_{\beta+1}\subseteq G$. 
Therefore, $G$ is a UCVL in $U$ containing $E$, and hence $F\subseteq G$ due to
Lemma \ref{Rem1}.

\smallskip
Let $V$ be a UCVL of $U$ containing $E$, so $F_1=E\subseteq V$. 
Let $\gamma$ be an ordinal such that $F_{\beta}\subseteq V$ for all $\beta<\gamma$.
If $\gamma$ is a limit ordinal then
$F_\gamma=\bigcup_{\beta<\gamma}F_{\beta}\subseteq V$.
If $\gamma=\beta+1$ then, for each $x\in F_\gamma$,
there exists a sequence $x_n$ in $F_\beta\subseteq V$
with $x_n \ruc x(x_1)$. As $V$ is UCVL then $x\in V$,
and since $x\in F_\gamma$ is arbitrary, $F_\gamma\subseteq V$.
In both cases, $G\subseteq\bigcup_{\gamma\in{\text{\rm Ord}}}F_{\gamma}\subseteq V$.
Since $V$ is arbitrary UCVL of $U$ containing $E$, then $G\subseteq F$
by Lemma \ref{Rem1}. In view of the already obtained inclusion
$F\subseteq G$, we conclude $F=G$.
\end{proof}

\begin{prop}[]\label{L2}
Let $E$ be a \text{\rm VSL} of a \text{\rm UCVL} $U$. 
Then the intersection $F$ of all \text{\rm UCVL}s of $U$ containing $E$
is an $\text{\rm r}$-completion of $E$.
\end{prop}

\begin{proof}
By Lemma \ref{L1}, $F$ is a UCVL that coincides with 
$\bigcup_{\gamma\in{\text{\rm Ord;$\gamma<\omega_1$}}}F_{\gamma}$.
In order to show that $F$ is an $\text{\rm r}$-completion of $E$,
let $G$ be a UCVL and let $T: E\to G$ be a Riesz 
homomorphism. It is enough to show that $T$ has
a unique extension to a Riesz homomorphism $S: F\to G$. 

\smallskip
First we prove that such an extension $S$,
if exists, is unique. Indeed assume $S_1: F\to G$ is another such an extension.
Clearly, $S$ agrees with $S_1$ on $F_1=E$. 
Assume $S$ agrees with $S_1$ on $F_{\beta}$ for all $\beta<\gamma$.
Let $x\in F_{\gamma}$.
If $\gamma$ is a limit ordinal then
$x\in F_{\gamma}=\bigcup_{\beta<\gamma}F_{\beta}$.
Then $x\in F_{\beta}$ for some $\beta<\gamma$, and hence
$Sx=S_1x$ by the assumption.
If $\gamma=\beta+1$ then there exists a sequence 
$x_n$ in $F_\beta$ with $x_n \ruc x(x_1)$. In other words, for each $k\in\mathbb{N}$, 
there exists $n_k$ with $|x_n-x|\le\frac{1}{k}|x_1|$ for all $n\ge n_k$.
Then 
$$
   |Sx_n-Sx|+|S_1x_n-S_1x|\le\frac{1}{k}(S|x_1|+S_1|x_1|)\ \ \ \ \ (\forall n\ge n_k).
   \eqno(2)
$$
Because the sequence $x_n$ is contained in $F_\beta$ then $Sx_n=S_1x_n$ 
for all $n\in\mathbb{N}$. Now, it follows from (2) that
$$
   Sx_n\ruc Sx(S|x_1|+S_1|x_1|) \ \ \ \text{\rm and} \ \ \ 
   S_1x_n\ruc S_1x(S|x_1|+S_1|x_1|)
$$ 
in $G$, and since $G$ is Archimedean, then $Sx=S_1x$ as desired.

\smallskip
In order to show the existence of the required extension $S: F\to G$,
we use Lemma \ref{L1} and apply the transfinite induction by $\gamma$.
Suppose the extensions $S_\alpha: F_\alpha\to G$ exist already
for all $\alpha<\gamma$. By the arguments as above,
they are unique for all $\alpha<\gamma$,
and hence $S_\alpha$ agree with $S_\beta$ on $F_{\min{(\alpha,\beta)}}$.
Let $x\in F_\gamma$.

\smallskip
(A) If $\gamma$ is a limit ordinal then $x\in F_\alpha$ 
for some $\alpha<\gamma$. We set $S_\gamma x:=S_\alpha x$.

\smallskip
(B) If $\gamma=\beta+1$, there exists a sequence 
$x_n$ in $F_\beta$ with $x_n \ruc x(x_1)$. Then, for each $k\in\mathbb{N}$, 
there exists $n_k$ with $|x_n-x|\le\frac{1}{k}|x_1|$ for all $n\ge n_k$, and hence 
$$
   |x_n-x_m|\le|x_n-x|+|x_m-x|\le\frac{2}{k}|x_1| \ \ \ \ (\forall n,m\ge n_k).
   \eqno(3)
$$
It follows from (3) that $|S_\beta x_n-S_\beta x_m|\le\frac{2}{k}S_\beta|x_1|\in G$ 
for all $n,m\ge n_k$.
So, the sequence $S_\beta x_n$ is $\text{\rm r}$-Cauchy in $G$, and hence there exists 
unique $y\in G$ with $S_\beta x_n\ruc y$ in $G$. 
It is easy to see that $y$ does not depend on the choice of the sequence 
$x_n$ in $F_\beta$ such that $x_n \ruc x(x_1)$. We set $S_\gamma x:=y$.

\smallskip
Finally, let $x\in F$. Then $x\in F_\gamma$ for some $\gamma$. 
We set $Sx:=S_\gamma x$. Since linear and lattice operations in VLs
are $\text{\rm r}$-continuous, the constructed extension $S: F\to G$ 
is a Riesz homomorphism. 
\end{proof}

\smallskip
Proposition \ref{L2} combined with Lemma \ref{princ-ideal} and with the
theorem of Kakutani and M. Krein -- S. Krein gives that each principal order ideal
of an $\text{\rm r}$-completion of an Archimedean VL is Riesz isomorphic
to $C(K)$ for some compact Hausdorff space $K$ (cf., \cite[Theorem 2]{Veks1}).


\subsection{Free \text{\rm r}-complete vector lattice over a set.}
Although the next definition does not, as such, 
appear in the literature, we think it belongs to folklore.

\begin{defi}[]\label{d1}
If $A$ is a non-empty set, then a {\em free UCVL over} $A$ 
{\em (}briefly, $\text{\rm FUCVL}(A)${\em )} is 
a pair $(F,f)$, where $F$ is a \text{\rm UCVL} and $f:A\to F$ is an embedding 
such that, for any \text{\rm UCVL} $E$ and for any embedding $q:A\to E$, 
there exists a unique Riesz homomorphism $T:F\to E$ satisfying $q=T\circ f$. 
\end{defi}

If exists, $\text{\rm FUCVL}(A)$ is an initial object in 
the category, whose objects are pairs $(E,q)$ with a UCVL $E$ and an embedding $q :A\to E$, 
and whose morphisms $q_2^1:(E_1, q_1)\to(E_2, q_2)$ are Riesz homomorphisms 
from $E_1$ to $E_2$ satisfying $q_2=q_2^1\circ q_1$. 
Thus, $\text{\rm FUCVL}(A)$ is defined similarly to $\text{\rm FVL}(A)$ in a proper subcategory. 
Routine arguments show that a free UCVL $(E,q)$ over $A$, if exists,
is unique up to a Riesz isomorphism.

\begin{thm}\label{maintt1}
Let $A$ be a non-empty set, $\text{\rm FVL}(A)$ be a \text{\rm VSL} 
of a \text{\rm UCVL} $U$, and $A\xrightarrow[]{i}\text{\rm FVL}(A)$ be 
the assigned embedding. Then$:$
\begin{enumerate}[$(i)$]
\item
the intersection $F$ of all \text{\rm UCVL}s of $U$ containing $\text{\rm FVL}(A)$  
together with the embedding $i$ is a $\text{\rm FUCVL}(A)$$;$
\item
$F=\bigcup\limits_{y\in\text{\rm FVL}(A)}U_y=
\bigcup\limits_{B\ \text{\rm is a nonempty finite 
subset of}\ A}U_{\sum\limits_{b\in B}|i(b)|}$.
\end{enumerate}
\end{thm}

\begin{proof}
$(i)$\
Note that $F$ is an $\text{\rm r}$-completion 
of $\text{\rm FVL}(A)$ due to Proposition \ref{L2}.

\smallskip
Let $E$ be a UCVL and let $q: A \to E$ be an embedding. 
By the definition of $\text{\rm FVL}(A)$, there exists 
a unique Riesz homomorphism $T_1:\text{\rm FVL}(A)\to E$
with $q = T_1\circ i$. Since $F$ is an $\text{\rm r}$-completion 
of $\text{\rm FVL}(A)$, there is 
unique Riesz homomorphism $T:F\to E$ satisfying $T_1 = T\circ J$,
where $J$ is the inclusion of $\text{\rm FVL}(A)$ into $F$, and hence
$q = T_1\circ i=T\circ J\circ i=T\circ i$.
By Definition \ref{d1}, the pair $(F,i)$ is a $\text{\rm FUCVL}(A)$.


\smallskip
$(ii)$\
It follows directly from Lemma \ref{princ-ideal}.
\end{proof}

Let $A$ be a non-empty set. Recall that $\text{\rm FVL}(A)$ is the VSL of 
UCVL $\R^{\R^A}$ generated by the evaluation functionals $\delta_a$ 
on $\R^A$, $\delta_a(\xi)=\xi(a)$ \cite[Theorem 2.4]{Baker}. The next 
representation of $\text{\rm FUCVL}(A)$ follows from Theorem \ref{maintt1}.

\begin{cor}\label{cortt1}
Let $A$ be a non-empty set and let $E$ be the \text{\rm VSL} of a \text{\rm UCVL} 
$\R^{\R^A}$ generated by evaluation functionals $\delta_a$, $a\in A$
with the embedding $A\xrightarrow[]{i}\R^{\R^A}$, $i(a):=\delta_a$. 
Then, the intersection $F$ of all \text{\rm UCVL}s of $\R^{\R^A}$ containing 
$i(A)$ together with embedding $i$ is a $\text{\rm FUCVL}(A)$.
\end{cor}


\subsection{Restrictions and projections.}
Let $A$ be a nonempty set. Following the tradition, for $\emptyset\ne B\subseteq A$, we identify 
$\R^{\R^B}$ with a VSL of  $\R^{\R^A}$
by assigning $\xi \in\R^{\R^B}$ with $\hat{\xi} \in\R^{\R^A}$ 
such that $\hat{\xi} (f) = \xi(f|_B)$ (see \cite{PW}).
By \cite[Proposition~3.5(2)]{PW}, there exists a unique Riesz 
homomorphism projection $P_B$ of $\text{\rm FVL}(A)$  onto $\text{\rm FVL}(B)$ satisfying
$$
   P_B(\delta_a)=
   \begin{cases}
   \delta_a  &  \text{if}  \ \  a\in B \\
   0  & \text{if}  \ \  a\in A\setminus B \, .
   \end{cases}
$$
In particular, 
$
   \text{\rm FVL}(A)= \bigcup\limits_{\emptyset\ne B\in{\cal P}_{fin}(A)}
   \text{\rm FVL}(B)\subseteq\R^{\R^A} 
$,
where ${\cal P}_{fin}(A)$ is the set of all finite subsets of $A$ (see \cite[Proposition 3.7]{PW}). 

\smallskip
Denote by $H(\R^A)$ (by $H(\Delta_A)$) the space of all positively 
homogeneous real-valued functions on $\R^A$ (on $\Delta_A:=[-1,1]^A$) 
which are continuous in the product topology of $\R^A$ (of $\Delta_A$). 
Clearly, $H(\Delta_A)$ is a VSL of $C(\Delta_A)$, closed in $\|.\|_\infty$. 
Thus, $H(\Delta_A)$ is a Banach lattice, and hence is a UCVL. 
It is well known that $\text{\rm FVL}(A)$ may be identified with a VSL of $H(\R^A)$ 
and hence with a VSL of $H(\Delta_A)$ in view of \cite[Lemma~5.1]{PW}. 
By \cite[Corollary~5.7]{PW}, $\text{\rm FBL}(A)$ is embedded into $H(\Delta_A)$ as 
a VSL $J(\text{\rm FBL}(A))$ and even as an order ideal, as mentioned in \cite{PW}. 
Furthermore,  $J(\text{\rm FBL}(A))=H(\Delta_A)$ iff $A$ is finite and, in this case,  
$\text{\rm FBL}(A)$ is isomorphic to $H(\Delta_A)$ under 
the supremum norm \cite[Theorem~8.2]{PW}.

\begin{thm}\label{l2}
Let $B$ be a non-empty finite set. Then $\text{\rm FUCVL}(B)$ is Riesz isomorphic 
to $\text{\rm FBL}(B)$, $H(\Delta_B)$, and to $H(\R^B)$. 
\end{thm}

\begin{proof}
As $\text{\rm FVL}(B)$ and $\text{\rm FBL}(B)$ possess (the same) strong 
order unit (cf. \cite[Proposition~5.3 and Theorem~8.2]{PW}) and $\text{\rm FBL}(B)$ 
is Riesz isomorphic to the completion of $\text{\rm FVL}(B)$ under its strong unit norm, 
$\text{\rm FBL}(B)$ is the $\text{\rm r}$-closure of $\text{\rm FVL}(B)$,
and hence $\text{\rm FUCVL}(B)$ is Riesz isomorphic to $\text{\rm FBL}(B)$.
The fact that $\text{\rm FBL}(B)$ 
and $H(\Delta_B)$ are Riesz isomorphic is the result of \cite[Proposition~5.3]{PW}. 
As $B$ is finite, the restriction map $R: H(\R^B) \to H(\Delta_B)$ is 
surjective and hence is a Riesz isomorphism (e.g. by \cite[Lemma~5.1]{PW}).
\end{proof}

For the convenience, we include the proof of the following certainly well known proposition.

\begin{prop}\label{p2}
Let $A\ne\emptyset$, and let $x\in\text{\rm FBL}(A)$. 
Then  there exists a sequence $x_n$ in $\text{\rm FVL}(A)$ 
that $\text{\rm r}$-converges to $x$ with a regulator $u\in\text{\rm FBL}(A)$.
\end{prop}

\begin{proof}
Since  $x\in\text{\rm FBL}(A)$, then there exists a sequence $x_n$ in $\text{\rm FVL}(A)$ 
satisfying $\|x_n - x\|_F \le \frac{1}{n2^n}$ for all $n\in \N$, where $\|.\|_F$ is 
the extension to $\text{\rm FBL}(A)$ of a free lattice norm on $\text{\rm FVL}(A)$. 
Let $u=\|.\|_F\text{-}\sum_{n=1}^{\infty}n |x_n -x|\in\text{\rm FBL}(A)$. 
Then $|x_n -x| \le\frac{1}{n}u$  for all $n\in \N$
and hence $x_n \ruc x(u)$.
\end{proof}

Since any Banach lattice is UCVL, it follows from Proposition \ref{p2} 
that $\text{\rm FBL}(A)$ contains $\text{\rm r}$-completion of $\text{\rm FVL}(A)$. 
We show that $\text{\rm FUCVL}(A)$ is a proper VSL 
of $\text{\rm FBL}(A)$ unless $A$ is finite. 

\begin{prop}\label{p3}
Let $A\ne\emptyset$. 
If a sequence $g_n \in H(\R^A)$  $\text{\rm r}$-converges with a regulator  
$u \in H(\R^A)$ to some  $g \in \R^{\R^A}$ then $g \in H(\R^A)$. 
\end{prop}

\begin{proof}
Without lost of generality, we may assume that $|g_n - g| \le \frac{1}{n} u$ for all $n\in \N$. 
Clearly, $g$ is positively homogeneous. It remains to show the continuity of $g$. 
Let $\eta_\alpha \to \eta$  in the product topology $\tau$ of $\R^A$. 
Then, for some $M$, there is $\alpha_0$ with $u(\eta_\alpha) \le M$  
for all $\alpha \ge \alpha_0$. Hence
$$
  |g(\eta_\alpha) - g(\eta)| \le |g(\eta_\alpha) -g_n(\eta_\alpha) |+ 
  |g_n(\eta_\alpha) -g_n(\eta) |+ |g_n(\eta) - g(\eta)|\le 
$$
$$
   |g_n(\eta_\alpha) -g_n(\eta) |+\frac{2M}{n}  
  \quad\quad  (\forall n\in \N) \quad (\forall \alpha \ge \alpha_0).
$$
Since $g_n$ is $\tau$-continuous, 
there exists $\alpha_n$ with $|g_n(\eta_\alpha) -g_n(\eta) |\le \frac{M}{n}$ 
for $\alpha \ge \alpha_n$. 
It follows that $|g(\eta_\alpha) - g(\eta)| \le \frac{3M}{n}$, 
if $\alpha \ge \alpha_0, \alpha_n$. Thus $g$ is $\tau$-continuous and hence $g \in H(\R^A)$.
\end{proof}

\begin{cor}\label{new}
$H(\R^A)$ is an $\text{\rm r}$-closed \text{\rm VSL} of $\R^{\R^A}$ 
with regulators from $H(\R^A)$.
\end{cor}

Considering $\text{\rm FVL}(A)$ as a VSL of $H(\Delta_A)$, 
we obtain the following result.

\begin{cor}\label{duble}
Let $R: H(\R^A)\to H(\Delta_A)$ be the restriction map.
If a sequence $g_n$ in $\text{\rm FVL}(A)$  $\text{\rm r}$-converges 
to some $g\in H(\Delta_A)$ with a regulator $u\in R(H(\R^A))$ then 
$g \in R(H(\R^A))$,  
\end{cor}

It is worth to compare Corollary \ref{duble} with  \cite[Example~5.2]{PW}. 

\subsection{Applications.}
The following theorem is a refinement of Theorem~\ref{l2}.

\begin{thm}\label{T1-1}
Let $A$ be a non-empty set. The intersection $F$ of all \text{\rm r}-closed 
\text{\rm VSL}s of $H(\R^A)$ containing $\text{\rm FVL}(A)$ is a $\text{\rm FUCVL}(A)$.
\end{thm}

\begin{proof}
By Corollary~\ref{new}, $H(\R^A)$ is a VSL of $\R^{\R^A}$
\text{\rm r}-closed under the regulators of $H(\R^A)$.
Since $\R^{\R^A}$ is a UCVL then $H(\R^A)$ is a UCVL.
As $\text{\rm FVL}(A)$ is a VSL of $H(\R^A)$, then
$\text{\rm FUCVL}(A)$ is an \text{\rm r}-closed in $H(\R^A)$ VSL of $H(\R^A)$.
\end{proof}

\begin{cor}\label{CorB}
Let $A$ be a non-empty set. Then 
$$
  \bigcup\limits_{\emptyset\ne B\in{\cal P}_{fin}(A)}\text{\rm FBL}(B)\subseteq 
  \text{\rm FUCVL}(A)\subseteq\text{\rm FBL}(A).
  \eqno(4)
$$
Furthermore, both inclusions are proper unless $A$ is finite. 
\end{cor}

\begin{proof}
Formula (4) follows from Lemma \ref{princ-ideal},
\cite[Theorem 8.2]{PW}, and Theorem \ref{l2} 
since every \text{\rm FBL} is an $\text{\rm UCVL}$. 
Let $A$ be infinite. The inclusion $\text{\rm FUCVL}(A)\subseteq\text{\rm FBL}(A)$ 
is proper since e.g. 
$g: = \sum_{k=1}^{\infty}2^{-k}\delta_{a_k}\in\text{\rm FBL}(A)$ for distinct $a_k \in A$
yet $g \not\in H(\R^A)$ (cf. \cite[Example~5.2]{PW}).
The inclusion $\bigcup\limits_{\emptyset\ne B\in{\cal P}_{fin}(A)}
\text{\rm FBL}(B)\subseteq\text{\rm FUCVL(A)}$
is proper since $z: = \sum\limits_{k=1}^\infty 2^{-k}|\delta_{a_k}|\wedge|\delta_{a_1}| 
\in\text{\rm FUCVL}(A)$ for distinct $a_k\in A$ yet $z\not\in\text{\rm FBL}(B)$ 
for each finite nonempty $B\subseteq A$.
\end{proof}


\begin{thebibliography}{30}
{\tiny 
\bibitem{AEG}%%% OK
A. Ayd{\i}n, E. Emelyanov, S.G. Gorokhova, 
Full lattice convergence on Riesz spaces, 
Indag. Math. (N.S.) 32 (2021), 658--690.

\bibitem{Baker}%%%OK
K.A. Baker, 
Free vector lattices,
Canad. J. Math. 20 (1968), 58--66.

\bibitem{Birk}%%%OK
G. Birkhoff,  
Lattice, ordered groups,
Ann. Math. 2 (43) (1942), 298--331.

\bibitem{Blei}%%%OK
R.D. Bleier,  
Free vector lattices, 
Trans. Am. Math. Soc. 176 (1973), 73--87. 

\bibitem{Ku}%%% OK
A.G. Kusraev,
Dominated Operators,
Kluwer Academic Publishers, Dordrecht (2000)

\bibitem{PW}%%%OK
B. de Pagter, A.W. Wickstead,
Free and projective Banach lattices, 
Proc. Roy. Soc. Edinburgh Sect. A 145 (2015), 105--143. 

\bibitem{Tr}%%%OK
V.G. Troitsky, 
Simple constructions of FBL(A) and FBL[E], 
Positivity 23 (2019), 1173--1178.

\bibitem{Veks1}%%%OK
A.I. Veksler,
The concept of a linear lattice which is normal in itself, and certain applications 
of this concept to the theory of linear and linear normed lattices (in Russian),
Izv. Vyssh. Uchebn. Zaved., Mat. No.4(53) (1966), 13--22.

\bibitem{Veks2}%%%OK
A.I. Veksler,
A new construction of Dedekind completion 
of vector lattices and of l-groups with division (in Russian), 
Siberian Math. J. 10 (1969), 891--896. 

\bibitem{Wein}%%%OK
E.C. Weinberg, 
Free lattice-ordered abelian groups,
Math. Annalen 151 (1963), 187--199.}
\end{thebibliography}

\end{document}
